All detector scores are converted to a higher-is-ID-like convention before
computing AUROC, AUPR-IN, and FPR95. This appendix lists the score definitions
used in the experiments and records whether each score is a faithful
reproduction of an existing method or a diagnostic variant used to probe a
geometry channel.

\subsection{Logit-Based Scores}
\label{app:logit-based-scores}

Logit-based detectors read output confidence, logit scale, and margin structure.
The maximum softmax probability score is
\[
    S_{\mathrm{MSP}}(x)
    =
    \max_c p_\theta(y=c\mid x).
\]
The Energy score is written in the higher-is-ID-like direction used in this
paper:
\[
    S_{\mathrm{Energy}}(x)
    =
    T\log\sum_{c=1}^{K}\exp(z_c(x)/T).
\]
The MaxLogit score is
\[
    S_{\mathrm{MaxLogit}}(x)
    =
    \max_c z_c(x).
\]
When negative entropy is reported, it is also oriented so larger values are more
ID-like.

\subsection{Distance-Based Feature Scores}
\label{app:distance-based-feature-scores}

The Mahalanobis detector reads class means and tied covariance through a
precision-weighted distance:
\[
    S_{\mathrm{Maha}}(x)
    =
    -\min_c
    (f_\theta(x)-\mu_c)^\top
    \Sigma_W^{-1}
    (f_\theta(x)-\mu_c).
\]
The negative sign enforces the higher-is-ID-like score convention. A
Mahalanobis-L2 variant applies the same score after normalizing features to unit
norm.

kNN detectors read local neighborhood distance in the ID training feature bank:
\[
    S_{\mathrm{kNN}}(x)
    =
    -d_k(f_\theta(x),\mathcal{F}_{\mathrm{tr}}),
\]
where \(d_k\) is the \(k\)-th nearest-neighbor distance. L2-kNN applies the same
score after feature normalization, allowing us to test whether the
local-distance detector depends on radial scale.

\subsection{Gaussian Feature-Density Scores}
\label{app:gaussian-feature-density-scores}

Gaussian feature-density detectors score a feature by a
class-conditional mixture likelihood:
\[
    S_{\mathrm{GMM}}(x)
    =
    \log \sum_{c=1}^{K}
    \pi_c
    \mathcal{N}(f_\theta(x)\mid \mu_c,\Sigma_c).
\]
We distinguish tied, diagonal, and shrinkage covariance variants as diagnostic
controls for covariance estimation and covariance conditioning.

We use Gaussian feature-density score to refer to the density scoring component
inspired by DDU. We do not claim to reproduce
the full original DDU training recipe unless spectral normalization and the
associated training protocol are explicitly used. This distinction is important
because spectral normalization is itself a training-time geometric regularizer
and can modify the optimizer-induced geometry that this paper aims to measure.

\subsection{Prototype, Angular, and Subspace Scores}
\label{app:prototype-angular-subspace-scores}

The reported results include NCC, prototype-cosine, and ViM diagnostics.
These are useful representation diagnostics but should not be relabeled as CTM
or NeCo.

CTM uses cosine similarity between features and classifier weights:
\[
    S_{\mathrm{CTM}}(x)
    =
    \max_c \cos(w_c,f_\theta(x))
    =
    \max_c
    \frac{w_c^\top f_\theta(x)}
    {\|w_c\|_2\|f_\theta(x)\|_2}.
\]
CTM reads angular alignment rather than feature norm and is directly connected
to NC3 self-duality between classifier weights and class means. In this
draft, CTM is a planned probe rather than an available main result.

NeCo uses the relationship between test features and the principal ID subspace.
With \(P\in\mathbb{R}^{d\times r}\) denoting the matrix of top ID principal
directions, the normalized projection score is
\[
    S_{\mathrm{NeCo}}(x)
    =
    \frac{\|P^\top f_\theta(x)\|_2}
    {\|f_\theta(x)\|_2}.
\]
When using NeCo, we distinguish the official score and fixed PCA configuration
from project-specific simplified variants. Temporary simplified scores are not
used as main results unless their formula, dimension, direction, and
normalization are fixed.

\subsection{Score Direction and Aggregation Convention}
\label{app:score-direction-aggregation}

All scores are oriented so larger values indicate more ID-like inputs before
AUROC, AUPR-IN, and FPR95 are computed. Aggregated near/far tables average over
the OOD datasets in the corresponding group after applying this score direction
convention.
